\section{Why large cyclic cores are the real next boundary}
\label{sec:beyond-bounded-blocks}

The fixed-face theorem and the response corollaries eliminate several false
blockers.  Restricted PageRank systems never have condition number worse than
$1/\alpha$; long spiders do not create an $\alpha^{-2}$ spectrum; and tree
support discovery can be performed without repeatedly solving every prefix.
What remains is the interaction of signed acceleration, changing faces, and
many coupled boundary demands inside a large biconnected core.

\subsection{A global relaxation parameter already handles all bipartite faces}

One might worry that every face change requires recomputing
$\sigma_U$ and $\omega_U^\star$.  A conservative global choice avoids that
particular cost.  Put
\begin{equation}
 t_\alpha:=\sqrt{1-c_\alpha^2}
 =\frac{2\sqrt\alpha}{1+\alpha},
 \qquad
 \omega_\alpha:=\frac2{1+t_\alpha},
 \qquad
 \zeta_\alpha:=\omega_\alpha-1
 =\left(\frac{1-\sqrt\alpha}{1+\sqrt\alpha}\right)^2.
 \label{eq:global-bipartite-omega}
\end{equation}

\begin{lemma}[One parameter for every fixed bipartite principal face]
\label{lem:one-omega-all-bipartite-faces}
For every bipartite principal face $U$, red--black SOR with
$\omega=\omega_\alpha$ has spectral radius at most $\zeta_\alpha$.  The
finite power bound
\begin{equation}
 \|T_U^k\|_2\leq6k\zeta_\alpha^{k-1}
 \label{eq:global-face-power-bound}
\end{equation}
holds without computing a face-specific eigenvalue.
\end{lemma}

\begin{proof}
Repeat the singular-mode proof of \Cref{thm:fixed-bipartite-sor} with the
upper bound $c_\alpha\sigma\leq c_\alpha$.  The identity
$\omega_\alpha^2c_\alpha^2=4\zeta_\alpha$ places every modal trace in
$[-2\zeta_\alpha,2\zeta_\alpha]$, with determinant
$\zeta_\alpha^2$.  The Cayley--Hamilton argument is unchanged.
\end{proof}

Thus spectral tuning is not the main obstacle on an evolving bipartite
support.  The difficulty is that $T_U$ acts on a different state space after
an admission, while the new restricted optimum has moved.

\subsection{The exact face shock at a settled checkpoint}

The simplest face-change identity shows what must be paid.  Fix a load
$\ell$; for RPPR on a positive face it is
$\ell=b-\alpha\rho D^{1/2}\mathbf1$.  Let $U'=U\mathbin{\dot\cup}E$, let
$z_U=Q_{UU}^{-1}\ell_U$, and zero-pad $z$ to $U'$.  Define the new-coordinate
demand
\begin{equation}
 g_E:=\ell_E-Q_{EU}z_U.
 \label{eq:new-face-demand}
\end{equation}

\begin{lemma}[Settled face-shock identity]
\label{lem:settled-face-shock}
If $z'=Q_{U'U'}^{-1}\ell_{U'}$, then
\begin{align}
 z'-z&=Q_{U'U'}^{-1}
       \begin{bmatrix}0\\g_E\end{bmatrix},
 \label{eq:face-shift-resolvent}\\
 \|z'-z\|_{Q_{U'U'}}^2
 &=g_E^T(z'_E-z_E)=g_E^Tz'_E.
 \label{eq:face-shock-energy}
\end{align}
If $g_E\geq0$, then $z'-z\geq0$.
\end{lemma}

\begin{proof}
The old equations are exact on $U$ and $z_E=0$, so
$\ell_{U'}-Q_{U'U'}z=(0,g_E)^T$.  Apply the inverse to obtain
\eqref{eq:face-shift-resolvent}; multiply by $(z'-z)^TQ_{U'U'}$ to obtain
\eqref{eq:face-shock-energy}.  The final assertion follows because a
principal SPD Stieltjes matrix has a nonnegative inverse.
\end{proof}

This identity supports an energy ledger when the old face is exactly settled.
But settling every old face before admitting one more coordinate recreates
the repeated-prefix tax.  A continuing accelerated method is generally not
at $z_U$ when $E$ arrives.  Its enlarged residual then has both old-face and
new-face components, and signed momentum may amplify or cancel them.  The
missing theorem must analyze that actual state, not repeatedly replace it by
the exact restricted optimum.

\subsection{Four distinct costs on a changing cyclic face}

A valid broad-graph proof must control all of the following.

\begin{enumerate}[leftmargin=*]
\item \textbf{Discovery.}  Which exterior coordinates violate an RPPR KKT
      inequality?  Reading every boundary row after every change can be
      quadratic in the eventual support.
\item \textbf{Numerical continuation.}  Does carried signed state contract
      over a window after zero-padding new coordinates, or must it be reset?
      A reset after every singleton admission loses acceleration.
\item \textbf{Representation.}  One local Schur update can alter many
      boundary demands.  Explicitly rewriting all of them is a cost even if
      the underlying matrix operation is algebraically low rank.
\item \textbf{Certification.}  A sufficient residual envelope may be much
      stronger than semantic error, as the signed star already shows.  The
      final test must not silently restore an avoidable logarithm.
\end{enumerate}

These costs explain why neither of the following implications is valid:
\begin{equation}
 \begin{aligned}
 \text{small final support}
 &\not\Longrightarrow
 \text{local accelerated trajectory},\\
 \text{fast fixed-face rate}
 &\not\Longrightarrow
 \text{fast face discovery}.
 \end{aligned}
 \label{eq:two-invalid-implications}
\end{equation}
The FISTA leaf-star refutes the first for one standard momentum rule.  The
literal endpoint-path restart examples refute the second for one safe
active-set implementation.

\subsection{The next positive structural classes}

The following classes form a sensible progression beyond bounded blocks.

\begin{enumerate}[leftmargin=*]
\item \textbf{Bipartite cores with a charged incremental gate.}  The global
      parameter \eqref{eq:global-bipartite-omega} settles spectral tuning.
      The only new theorem should concern face shocks, KKT reporting, and
      state transport.
\item \textbf{Nearly bipartite cores.}  Choose a partition and write
      $S=S_{\rm cross}+S_{\rm same}$.  A useful conditional target is to
      absorb $S_{\rm same}$ as a perturbation when its operator norm or its
      locally charged energy is $O(\sqrt\alpha)$.  A statement based only on
      the number of same-color edges is too crude for high-degree graphs.
\item \textbf{Graphs with small dynamic separators or response rank.}
      Maintain a low-dimensional Schur response at the separator and run SOR
      only on the unsettled side.  Every separator rebuild and response query
      must be included; low final rank alone is not a maintenance theorem.
\item \textbf{Arbitrary large cores.}  Exact boundary refresh is likely too
      strong.  The weakest sufficient interface is an output-sensitive
      reporter that identifies only KKT demands crossing a finite accuracy
      band, with missed light changes controlled by an energy packing bound.
\end{enumerate}

The last item is compatible with the existing certified-response direction:
one need not maintain every exact zero sign.  It remains open to combine such
a reporter with a signed iterative state while retaining the product scale.

\subsection{Mandatory counterexample suite}

Every proposed general theorem should be tested against a different failure
mechanism on each of the following families.

\begin{center}
\begin{tabular}{p{0.24\textwidth}p{0.31\textwidth}p{0.34\textwidth}}
\toprule
Family & What it stresses & A proof must demonstrate \\
\midrule
Center-star & Slow monotone transport and semantic cancellation & Signed
progress without replacing semantic error by a stronger residual test \\
Leaf-star & Momentum overshoot onto a high-degree hub & Confinement,
safeguarding, or an explicit charge for the transient hub scan \\
Endpoint path & Repeated growing-prefix work & Persistent state or a
geometric/windowed amortization \\
Asymmetric T tree & Debt surviving an admission & A multistage rather than
one-step reset potential \\
Comb or high-rank boundary & Many distinct response directions & Compressed
reporting, not an uncharged dense boundary vector \\
Bounded-degree expander & Large fronts and no small articulation & A genuine
large-core mechanism, or an explicit structural exclusion \\
\bottomrule
\end{tabular}
\end{center}

A failure on one named recurrence should remain scoped to that recurrence.
Conversely, passing finite examples is not a proof of the required windowed
or graph-uniform inequality.
